\section{Smoothing two contact lines}
\label{sm:section}

Choose the two free contact lines
$f_i:\PP^1\to X$ supplied by Proposition~\ref{lines:suitable-pair}.
Write $f_1(0)=f_2(0)=x$, and choose affine parameters $u,v$ on
their sources. In target coordinates centered at $x$, and a local
frame of $L$ representing $\theta$ by $\alpha$, put
\begin{equation}\label{sm:pair}
a=f_1'(0),\qquad b=f_2'(0),\qquad
d\alpha_x(a,b)\ne0,\qquad
(P_{f_1})_0+(P_{f_2})_0=F_x.
\end{equation}
Here $P_{f_i}$ is the positive subbundle of $f_i^*T_X$.
We will smooth this pair while allowing its two outer values to
move. The resulting family will contain a map of contact degree two
that is \emph{very free}: every splitting summand of its pulled-back
tangent bundle has positive degree.

\subsection{The source family and the evaluation cokernel}

Use homogeneous coordinates $[U_0:U_1]$, $[V_0:V_1]$, with
$u=U_1/U_0$ and $v=V_1/V_0$. Consider
\begin{equation}\label{sm:source}
\mathscr C=
\{U_1V_1=sU_0V_0\}
\subset\PP^1\times\PP^1\times\mathbb A^1_s,
\qquad \pi:\mathscr C\longrightarrow\mathbb A^1_s.
\end{equation}
The affine equation $uv=s$ and the charts at infinity show
that this is a projective flat family of nodal genus-zero curves.
Its central fiber $C=C_1\cup C_2$ consists of the two axes
$C_1=\{v=0\}$ and $C_2=\{u=0\}$, meeting at $o=(0,0)$.
For $s\ne0$ the fiber is smooth, and projection to the first
factor identifies it with $\PP^1$. There are disjoint sections
\[
p(s)=(\infty,0),\qquad q(s)=(0,\infty),
\]
both in the relative smooth locus. The maps $f_1$ and $f_2$
agree at $o$, so glue to $f_0:C\to X$.
Set $E=f_0^*T_X$ and $E_i=f_i^*T_X$.

We recall precisely the deformation statement needed here.

\begin{lemma}[Relative deformation of maps]\label{sm:relative-hom}
Let $S$ be a smooth irreducible complex curve, let
$\mathscr Y\to S$ be a projective flat family of nodal curves,
and let $Z$ be a smooth projective variety. The relative scheme
of maps
\[
\operatorname{Hom}_S(\mathscr Y,Z\times S)\longrightarrow S
\]
is locally of finite type. At a map $f:Y_{s_0}\to Z$ its
relative tangent space is $H^0(Y_{s_0},f^*T_Z)$.
If $H^1(Y_{s_0},f^*T_Z)=0$, the displayed morphism is smooth
at $[f]$. In particular its total space is smooth there and
its differential onto $T_{s_0}S$ is surjective.
\end{lemma}

\begin{proof}
This is the relative-Hom smoothness Theorem of
\cite[Theorem~15]{AraujoKollar2003}, specialized to a constant
smooth target. That Theorem requires a proper flat source,
a smooth relative target, and a central source without embedded
points. Our projective flat source and its nodal central fiber
satisfy these hypotheses; the source need not be smooth.
Representability follows from the graph construction.
For completeness, local lifts of a map over a small square-zero
extension exist by smoothness of the target. Their differences
are sections of $f^*T_Z$ tensored with the extension ideal, and
the obstruction to gluing them lies in its first cohomology.
Vanishing of $H^1(f^*T_Z)$ gives formal smoothness, hence
smoothness for this locally finite-type scheme of maps; the
same description gives the asserted tangent space.
This concerns the full relative scheme of maps. It does not
require choosing an extension over the original base without
an \emph{\'etale} base change; compare
\cite[Warning~16 and Application~17]{AraujoKollar2003}.
\end{proof}

The normalization sequence is
\begin{equation}\label{sm:normalization}
0\longrightarrow E
\longrightarrow (i_1)_*E_1\oplus(i_2)_*E_2
\xrightarrow{\ \operatorname{ev}_o-\operatorname{ev}_o\ }
T_xX\otimes\OO_o\longrightarrow0,
\end{equation}
where $i_i:C_i\to C$ are the inclusions.
Freeness gives $H^1(C_i,E_i)=0$ and surjectivity of evaluation
at $o$. Thus
\begin{equation}\label{sm:unconstrained-vanishing}
H^1(C,E)=0.
\end{equation}
Let $\mathscr H$ denote the relative scheme of maps for
\eqref{sm:source}, near $[f_0]$, and let
$\sigma:\mathscr H\to\mathbb A^1$ be its projection.
Lemma~\ref{sm:relative-hom} makes $\sigma$ smooth at $[f_0]$.
Evaluation at the two sections defines
\[
\operatorname{Ev}:\mathscr H\longrightarrow X\times X,
\qquad [f:\mathscr C_s\to X]\longmapsto(f(p(s)),f(q(s))).
\]
Put $y_1=f_1(\infty)$ and $y_2=f_2(\infty)$.

\begin{lemma}\label{sm:vertical-cokernel}
The cokernel of evaluation on the vertical tangent space is
canonically
\begin{equation}\label{sm:cokernel}
\coker\bigl[
H^0(C,E)\longrightarrow T_{y_1}X\oplus T_{y_2}X
\bigr]
\simeq H^1(C,E(-p-q))
\simeq T_xX/F_x.
\end{equation}
In particular it has dimension one.
\end{lemma}

\begin{proof}
Here $p,q$ also denote the corresponding smooth points of $C$.
Twisting the normalization sequence gives
\[
0\longrightarrow E(-p-q)
\longrightarrow(i_1)_*E_1(-p)\oplus(i_2)_*E_2(-q)
\longrightarrow T_xX\otimes\OO_o\longrightarrow0.
\]
All splitting degrees of $E_i$ are nonnegative. Consequently
$H^1(E_1(-p))=H^1(E_2(-q))=0$, and the images at $o$ of
their spaces of sections are exactly $(P_{f_1})_0$ and
$(P_{f_2})_0$: the trivial summands contribute no sections after
the twist, and each positive summand remains generated at $o$.
Equation~\eqref{sm:pair} identifies the resulting cokernel
with $T_xX/F_x$. Finally the sequence
\[
0\longrightarrow E(-p-q)\longrightarrow E
\longrightarrow E|_p\oplus E|_q\longrightarrow0
\]
and \eqref{sm:unconstrained-vanishing} identify the same group
with the evaluation cokernel in \eqref{sm:cokernel}.
\end{proof}

\subsection{The obstruction with both ends fixed}

It remains to show that smoothing supplies this missing evaluation
direction; we do so by ruling out a first-order smoothing that fixes
both endpoint values.

\begin{lemma}\label{sm:fixed-ends}
There is no tangent vector $\zeta\in T_{[f_0]}\mathscr H$
such that
\[
d\sigma(\zeta)=1,\qquad d\operatorname{Ev}(\zeta)=0.
\]
\end{lemma}

\begin{proof}
Such a vector would give a map to $X$ from the pullback of
\eqref{sm:source} under $s=\epsilon$, where
$\epsilon^2=0$, reducing to $f_0$ and fixing the two endpoint
values to first order. Near the node its source has equation
\[
uv=\epsilon.
\]
On $C_1\setminus\{o\}$ this first-order source is trivialized
by $u$, including at infinity: in the coordinate $w=1/u$ its
second coordinate is $v=\epsilon w$.
The infinitesimal difference from $f_1$ is therefore an
intrinsic section $A$ of $E_1$ on this punctured component.
Similarly there is a section $D$ of $E_2$ on
$C_2\setminus\{o\}$. The fixed endpoint conditions say
\begin{equation}\label{sm:endpoint-zeros}
A(\infty)=D(\infty)=0.
\end{equation}

Use the same target coordinates centered at $x$ on both
branches. Let $g_1(u),g_2(v)$ be the vector-valued coordinate
expressions of $f_1,f_2$; their constant terms are zero.
In the completed local ring
\[
\CC[[u,v,\epsilon]]/(uv-\epsilon,\epsilon^2),
\]
the hypothetical map has coordinate expression
\begin{equation}\label{sm:node-lift}
g_1(u)+g_2(v)+\epsilon R(u,v)
\end{equation}
for a vector of formal power series $R$.
Indeed its reduction modulo $\epsilon$ agrees with $g_1$
and $g_2$ on the two axes, so the remaining difference is
divisible by $\epsilon$.
Substitute $v=\epsilon/u$, respectively $u=\epsilon/v$,
in \eqref{sm:node-lift}. Since $g_1'(0)=a$, $g_2'(0)=b$,
and $\epsilon^2=0$, the Laurent expansions are
\begin{equation}\label{sm:shared-constant}
A(u)=\frac b u+c+O(u),\qquad
D(v)=\frac a v+c+O(v),\qquad c=R(0,0).
\end{equation}
The vector $c$ is independent of the choice of $R$. Indeed, the
local ring is $\CC[[u,v]]/((uv)^2)$, with $\epsilon=uv$ and
$\operatorname{Ann}(\epsilon)=(\epsilon)$, so an ambiguity in $R$
has zero constant term. Thus the regular constants agree.
Mixed monomials in $\epsilon R$ do not affect them, since
$\epsilon uv=\epsilon^2=0$.
This equality also survives a change of target coordinates:
if its derivative and Hessian at $x$ are $J$ and $B$, the two
constants become $Jc+B(a,b)$ and $Jc+B(b,a)$, which agree
by symmetry of $B$.

Consider the sections
\[
W_1=uA\quad\hbox{and}\quad W_2=vD.
\]
They extend over the node by \eqref{sm:shared-constant}.
At infinity the simple poles of $u$ and $v$ are canceled by
the zeros in \eqref{sm:endpoint-zeros}. Thus
$W_i\in H^0(C_i,E_i)$ are global sections. In our common
target coordinates their jets at the node are
\begin{equation}\label{sm:auxiliary-jets}
W_1(0)=b,\quad W_1'(0)=c,\qquad
W_2(0)=a,\quad W_2'(0)=c.
\end{equation}

Freeness makes the ordinary scheme of maps smooth at each
$f_i$, so $W_i$ is the velocity of a local deformation of
that degree-one map. Every member of each such deformation
is horizontal: its contact differential is a homomorphism
$\OO(2)\to\OO(1)$ and hence vanishes.
We use the infinitesimal horizontality calculation in the proof of
\cite[Proposition~3.1]{Kebekus2001}. For a variation of $f_i$ with
velocity $W_i$, differentiation of horizontality in the chosen frame
and target coordinates gives
\[
0=(\partial_{W_i}\alpha)(f_i')+\alpha(W_i').
\]
Here $\partial$ differentiates the coefficient functions of
$\alpha$. Evaluating at the node and using
\eqref{sm:auxiliary-jets} yields
\[
(\partial_b\alpha)_x(a)+\alpha_x(c)=0,
\qquad
(\partial_a\alpha)_x(b)+\alpha_x(c)=0.
\]
Subtracting gives $d\alpha_x(a,b)=0$, contradicting
\eqref{sm:pair}. Horizontality was used for the auxiliary
degree-one deformations, whose velocities are $W_1,W_2$.
\end{proof}

\subsection{A very free conic and its full neighborhood}

\begin{proposition}\label{sm:very-free}
There is a map $g:\PP^1\to X$ with $\deg g^*L=2$ and
\begin{equation}\label{sm:positive-splitting}
g^*T_X\simeq\OO(2)\oplus\OO(1)^{\oplus2n}.
\end{equation}
The full scheme $\Hom_2(\PP^1,X)$ of degree-two maps is smooth
at $g$, and evaluation at two fixed distinct parameters is
submersive there. Its irreducible component through $g$
therefore dominates $X\times X$ by two-point evaluation.
Moreover $g$ is an immersion, and strict positivity of the
tangent splitting persists on an open neighborhood of $g$.
\end{proposition}

\begin{proof}
Let $T=T_{[f_0]}\mathscr H$ and $T_0=\ker d\sigma$.
Lemma~\ref{sm:relative-hom} supplies $\zeta\in T$ with
$d\sigma(\zeta)=1$. By Lemma~\ref{sm:vertical-cokernel},
the subspace $V=d\operatorname{Ev}(T_0)$ has codimension one
in $T_{y_1}X\oplus T_{y_2}X$.
If $d\operatorname{Ev}(\zeta)$ belonged to $V$, subtracting
a vector of $T_0$ would contradict Lemma~\ref{sm:fixed-ends}.
Thus the image of $\zeta$ supplies the missing direction and
$d\operatorname{Ev}:T\to T_{y_1}X\oplus T_{y_2}X$ is
surjective. Since its source and target are smooth at the
points in question, $\operatorname{Ev}$ is smooth at $[f_0]$.

Shrink to a neighborhood $\mathscr U$ on which both
$\sigma$ and $\operatorname{Ev}$ are smooth. This neighborhood
meets $s\ne0$, since a smooth morphism is open.
The contact degree of maps in $\mathscr U$ is two: degree is
locally constant in a projective flat family carrying a line
bundle, and its value at $f_0$ is $1+1$.
Over $\mathbb G_m$ the source family is explicitly the product
$\PP^1\times\mathbb G_m$, with isomorphism
\begin{equation}\label{sm:trivialization}
([U_0:U_1],s)\longmapsto
\bigl([U_0:U_1],[U_1:sU_0],s\bigr).
\end{equation}
The sections $p,q$ correspond to the fixed parameters
$\infty,0$. Consequently the degree-two relative scheme of
maps on this locus is
\[
\Hom_2(\PP^1,X)\times\mathbb G_m,
\]
and $\operatorname{Ev}$ is $(h,s)\mapsto(h(\infty),h(0))$.
In particular its differential in the $s$ direction is zero.
Choose a point $(g,s)\in\mathscr U$ with $s\ne0$.
Smoothness of $\sigma$ makes the full scheme of maps smooth
at $g$, and surjectivity of $d\operatorname{Ev}$ implies
surjectivity of
\begin{equation}\label{sm:two-values}
H^0(\PP^1,g^*T_X)
\longrightarrow T_{g(\infty)}X\oplus T_{g(0)}X.
\end{equation}
Thus the base parameter contributes no artificial evaluation
direction.

Write $g^*T_X=\bigoplus_{j=1}^d\OO(a_j)$, where $d=2n+1$.
For a line bundle on $\PP^1$, evaluation of sections at two
distinct points is surjective exactly when its degree is at
least one. Equation~\eqref{sm:two-values} therefore gives
$a_j\geq1$ for every $j$. Also
\[
\sum_{j=1}^d a_j
=\deg g^*(-K_X)=2(n+1)=d+1,
\]
which proves \eqref{sm:positive-splitting}.
The unique local irreducible component of the smooth scheme
of maps at $g$ has a submersive two-point evaluation, hence
dominates $X\times X$. The nonconstant map $g$ has nonzero
differential in characteristic zero. In
\eqref{sm:positive-splitting} a homomorphism
$T_{\PP^1}=\OO(2)\to g^*T_X$ takes values in the unique
degree-two summand, and is a nonzero constant isomorphism
onto that summand. Thus $g$ is an immersion.
Finally strict positivity is equivalent to
$H^1(\PP^1,g^*T_X(-2))=0$, an open condition by
semicontinuity. This proves the neighborhood assertion.
\end{proof}

We record how this full neighborhood passes to the unmarked
stable-map space used below.

\begin{lemma}\label{sm:full-germ}
Let $\beta$ be the numerical curve class of $g$, specified by
its degree functional on $\Pic(X)$. The stack
$\overline{\mathcal M}_{0,0}(X,\beta)$ is smooth at $[g]$,
and has a unique irreducible component through that point.
If $\mathcal S$ is the normalization of that reduced
component, it agrees with the original stack on an open
neighborhood of $[g]$. After shrinking an open neighborhood
$\mathcal H$ of $g$ in $\Hom_2(\PP^1,X)$, the map forgetting
parametrization factors through $\mathcal S$, and the
universal stable-map family pulls back to the full family
$\mathcal H\times\PP^1\to X$.
\end{lemma}

\begin{proof}
Since $\Pic(X)=\ZZ[L]$, contact degree two determines $\beta$,
so the class-$\beta$ scheme of maps is $\Hom_2(\PP^1,X)$.
A nonconstant map from $\PP^1$ has finite
automorphism group and is stable without marks.
On the open locus of smooth sources the stable-map stack is
the quotient of the corresponding scheme of parametrized
maps by $\PGL_2$: smooth families of rational curves acquire
parametrizations locally in the smooth topology.
By Proposition~\ref{sm:very-free} and smoothness of $\PGL_2$,
this quotient is smooth near $[g]$. It is consequently normal
and has a unique local irreducible component there.
Normalization of the reduced component is therefore an
isomorphism on that neighborhood. Pulling it back to the
scheme of maps gives an open neighborhood $\mathcal H$ of
$g$, with its original universal family. The neighborhood
$\mathcal H$ is open in the full scheme $\Hom_2(\PP^1,X)$, so
its tangent space at each $h\in\mathcal H$ is still
$H^0(\PP^1,h^*T_X)$. This retains every infinitesimal variation
used in the first-jet argument below.
\end{proof}
